\documentclass[11pt,a4paper]{article}

% ---------- Packages ----------
\usepackage{fontspec}
\usepackage[margin=2.2cm]{geometry}
\usepackage{amsmath,amssymb}
\usepackage{booktabs,tabularx,array,longtable}
\usepackage{enumitem}
\usepackage[dvipsnames,table]{xcolor}
\usepackage{listings}
\usepackage{titlesec}
\usepackage{fancyhdr}
\usepackage{tcolorbox}
\usepackage{parskip}
\usepackage{hyperref}

% ---------- Colours ----------
\definecolor{flame}{HTML}{E4572E}
\definecolor{ember}{HTML}{F3A712}
\definecolor{ink}{HTML}{1F2933}
\definecolor{soft}{HTML}{FFF4EC}
\definecolor{codebg}{HTML}{F6F7F9}

\hypersetup{colorlinks=true,linkcolor=flame,urlcolor=flame,
  pdftitle={Flamme Build Plan},pdfauthor={Team Flamme}}

% ---------- Headings ----------
\titleformat{\section}{\Large\bfseries\color{flame}}{\thesection}{0.8em}{}[\vspace{-0.4em}{\color{ember}\rule{\linewidth}{0.8pt}}]
\titleformat{\subsection}{\large\bfseries\color{ink}}{\thesubsection}{0.7em}{}
\titlespacing*{\section}{0pt}{1.6em}{0.9em}
\titlespacing*{\subsection}{0pt}{1.2em}{0.5em}

% ---------- Header / footer ----------
\pagestyle{fancy}
\fancyhf{}
\fancyhead[L]{\small\color{gray}Flamme --- Build Plan}
\fancyhead[R]{\small\color{gray}Team Flamme}
\fancyfoot[C]{\small\thepage}
\fancyfoot[R]{\scriptsize\color{gray}Not a NASA product}
\renewcommand{\headrulewidth}{0.3pt}

% ---------- Lists, tables, code ----------
\setlist{itemsep=2pt,topsep=4pt}
\renewcommand{\arraystretch}{1.25}
\newcolumntype{L}{>{\raggedright\arraybackslash}X}
\lstset{basicstyle=\ttfamily\small,backgroundcolor=\color{codebg},frame=single,
  rulecolor=\color{gray!40},breaklines=true,columns=fixed,keepspaces=true,basewidth=0.52em,
  keywordstyle=\color{flame}\bfseries,commentstyle=\color{gray},xleftmargin=4pt,xrightmargin=4pt}

% ---------- Helpers ----------
\newcommand{\Otwo}{\ensuremath{\mathrm{O_2}}}
\newcommand{\code}[1]{\texttt{#1}}
\newtcolorbox{donebox}{colback=soft,colframe=ember,boxrule=0.6pt,arc=2pt,
  left=6pt,right=6pt,top=3pt,bottom=3pt}
\newcommand{\done}[1]{\begin{donebox}\textbf{Done when:} #1\end{donebox}}
\newtcolorbox{keybox}{colback=soft,colframe=flame,boxrule=0.8pt,arc=2pt,left=8pt,right=8pt}

\setlength{\emergencystretch}{3em}\tolerance=2000
\begin{document}

% ================= Title =================
\begin{titlepage}
\centering
\vspace*{3cm}
{\color{flame}\rule{\linewidth}{1.5pt}}\\[1em]
{\fontsize{34}{40}\selectfont\bfseries\color{ink} Flamme\par}
\vspace{0.4em}
{\LARGE Build Plan\par}
\vspace{0.8em}
{\large From prototype to evidence platform:\\ stack, architecture and step-by-step build\par}
{\color{flame}\rule{\linewidth}{1.5pt}}\\[2.5em]
{\large Team Flamme\par}
\vspace{0.3em}
NASA Space Apps Challenge $\cdot$ \emph{Flame in Freefall}\par
\vspace{0.3em}
October 2026\par
\vfill
\begin{keybox}
\small\textbf{Guiding rule.} Every value shown is either \textbf{Observed} (cited source row or page),
\textbf{Predicted} (named model version with uncertainty) or \textbf{Unknown}.
Nothing is ever silently filled in. Flamme is not a NASA product and must not be used for
material certification or operational safety decisions.
\end{keybox}
\end{titlepage}

\tableofcontents
\newpage

% ================= 0 =================
\section{What we are building}

The plan document defines six deliverables on top of the existing prototype.

\begin{tabularx}{\linewidth}{@{}c p{3.6cm} L@{}}
\toprule
\# & \textbf{Deliverable} & \textbf{What it means technically}\\
\midrule
1 & Evidence database & Every NASA experiment stored in the ten-condition schema, with a source link and a per-field ``not reported'' flag\\
2 & Search and compare & Filter by any condition (``PMMA, $<21\%$ \Otwo, opposed flow $<10$ cm/s''); side-by-side view in standard units\\
3 & Gap map & Coverage plot of tested versus untested condition combinations\\
4 & Data-grounded 3D lab & Replace the hand-tuned curves in \code{labModel.js} with relationships fitted to the database; show the nearest real experiments beside the simulation\\
5 & Reporting template & The ten-condition format published as a JSON Schema and a fillable template\\
6 & Outreach mode & Keep the 92-second guided Simulate sweep for classrooms and the public\\
\bottomrule
\end{tabularx}

\subsection*{Reused from the current repository (not rewritten)}
\begin{itemize}
  \item \code{src/Lab.jsx}, \code{src/labScene.js}, \code{src/labModel.js} --- the Three.js flame lab and 92-second timeline.
  \item \code{src/data.js}, \code{src/science.js} --- the field model (\code{source / unit / normalized / status}) and transparent matching. This field model becomes the database schema.
  \item \code{flame\_in\_freefall\_data/psi\_scraper.py}, \code{download\_ntrs.py} --- PSI and NTRS harvesting (24 investigations, about 3{,}000 table rows, 36 PDFs).
  \item \code{findings/} --- table-to-finding extraction and the PDF extraction utility.
  \item \code{parameter\_catalog/} --- the ten-condition reference.
\end{itemize}

% ================= 1 =================
\section{The stack}

\begin{tabularx}{\linewidth}{@{}p{3.2cm} p{4.6cm} L@{}}
\toprule
\textbf{Layer} & \textbf{Choice} & \textbf{Why}\\
\midrule
Frontend & React 19 + Vite 7 + TypeScript & Already used; TypeScript catches API/UI schema mismatches. Migrate gradually (\code{allowJs})\\
3D & Three.js (existing \code{labScene.js}) & Already built; no rewrite needed\\
Routing / fetching & TanStack Router + TanStack Query & Typed, shareable search URLs; caching\\
Lab state & Zustand & Small; readable from the Three.js render loop\\
Tables & TanStack Table & Sortable, filterable compare view\\
Charts & Observable Plot (or ECharts) & Heatmaps and scatter plots with little code\\
Icons / style & Lucide, existing \code{app.css} as CSS variables & Keep the current look\\
\midrule
API & Python 3.12 + FastAPI + Pydantic v2 & Pipeline and ML are already Python; Pydantic generates the JSON Schema for deliverable 5\\
ORM / migrations & SQLAlchemy 2 + Alembic & Versioned schema changes\\
Database & PostgreSQL 16 on Neon or Supabase & Relational, range queries, JSONB for raw source rows\\
Units & \code{pint} & Safe atm/psi/kPa, \%/mole-fraction, mm/cm, \textdegree C/K conversion\\
\midrule
Wrangling & pandas (or polars) + pandera & Per-investigation adapters with schema checks\\
PDF extraction & PyMuPDF + pdfplumber + LLM structured output & Pull conditions from NTRS reports with page and quote; human review required\\
ML & scikit-learn, statsmodels, optional PyMC & Small data: simple, interpretable, calibrated models\\
Tracking & MLflow or versioned JSON in \code{models/} & Every fit has a version the UI can cite\\
Data versioning & DVC (Drive / S3 / R2 remote) & 239 MB of raw data must not live in git\\
\midrule
Auth & Supabase Auth (or Clerk) & Public site is read-only; only reviewers log in\\
Tests & Vitest, Playwright, pytest & Unit, browser and pipeline coverage\\
CI & GitHub Actions & Lint, test, build and data validation on every PR\\
Hosting & Vercel (web), Render/Fly/Railway (API), Neon (DB) & Free tiers, push-to-deploy\\
\bottomrule
\end{tabularx}

\begin{keybox}
\textbf{Cheaper alternative for a hackathon window.} Skip the live API. Run the pipeline at build time,
export \code{experiments.json} and \code{models/*.json}, and query in the browser with
\textbf{DuckDB-WASM}. Same schema, same pipeline, no server. Move to FastAPI + Postgres when
the reviewer workflow (Step~4) is needed.
\end{keybox}

% ================= 2 =================
\section{Repository layout (target)}

\begin{lstlisting}
flamme/
|-- apps/
|   |-- web/                  React + Vite + TS (moved from src/)
|   |   |-- src/lab/          Lab.tsx, labScene.ts, labModel.ts
|   |   |-- src/search/       filter bar, results, compare view
|   |   |-- src/gaps/         gap map
|   |   `-- src/api/          generated types + TanStack Query hooks
|   `-- api/                  FastAPI
|       |-- flamme_api/
|       |   |-- models.py     SQLAlchemy tables
|       |   |-- schemas.py    Pydantic: source of truth for the 10 conditions
|       |   |-- routes/       experiments, search, gaps, models, review
|       |   `-- query.py      filter -> SQL
|       `-- alembic/
|-- pipeline/
|   |-- harvest/              psi_scraper.py, download_ntrs.py (moved)
|   |-- adapters/             one per investigation: saffire2.py, bass2.py ...
|   |-- extract_pdf/          NTRS report extraction (LLM + quote + page)
|   |-- normalize.py          units via pint, vocabularies, status flags
|   |-- validate.py           pandera checks
|   `-- load.py               upsert into Postgres
|-- ml/
|   |-- features.py
|   |-- fit_flammability.py
|   |-- fit_spread.py
|   |-- fit_suppression.py
|   `-- export.py             -> models/<name>@<version>.json for the lab
|-- schema/
|   `-- flamme-experiment.schema.json    (generated; deliverable 5)
|-- data/                     DVC-tracked raw + processed
`-- docs/
\end{lstlisting}

% ================= 3 =================
\section{The core data model}

\emph{Build this first --- everything else depends on it.} Every condition value is a
\textbf{Field}, exactly the shape already used in \code{src/data.js}:

\begin{lstlisting}[language=Python]
class Field(BaseModel):
    source: str | None          # verbatim: "~ 22.1", "037 mm", "Insignificant"
    value: float | str | None   # normalized, only when defensible
    unit: str | None            # standard unit (table below)
    status: Literal["reported", "approximate", "qualitative",
                    "ambiguous", "not_reported", "not_applicable", "derived"]
    note: str | None = None
\end{lstlisting}

\subsection{The ten conditions}

\begin{tabularx}{\linewidth}{@{}c p{3.4cm} L@{}}
\toprule
\# & \textbf{Condition} & \textbf{Fields (standard unit)}\\
\midrule
1 & Material / fuel & name, material class (cellulosic, PMMA, silicone, fabric composite, liquid droplet, gas, \dots), formulation, coating, substrate\\
2 & Geometry & length, width, thickness (mm); thermally thin/thick flag (derived)\\
3 & Gravity & regime (Earth, lunar, Martian, microgravity), $g$ level ($g_0$), platform (ISS, Cygnus, drop tower, parabolic)\\
4 & Oxygen & $X_{\mathrm{O_2}}$ (\% vol)\\
5 & Pressure & $P$ (kPa)\\
6 & Forced airflow & $u$ (cm/s)\\
7 & Flow direction & concurrent, opposed, cross, quiescent\\
8 & Ignition & igniter type, power (W), duration (s), position\\
9 & Initial temperature & $T_0$ (K)\\
10 & Suppression & agent, concentration (\% vol), flow rate, timing (s after ignition)\\
\bottomrule
\end{tabularx}

\medskip
\noindent\textbf{Derived column.} Oxygen partial pressure, because the plan states that it, not percentage alone,
drives burning:
\[
  p_{\mathrm{O_2}} \;=\; \dfrac{X_{\mathrm{O_2}}}{100}\, P
\]
It is stored with status \code{derived} and only when \emph{both} inputs are reported.

\subsection{Database tables}

\begin{tabularx}{\linewidth}{@{}p{3.4cm} L@{}}
\toprule
\textbf{Table} & \textbf{Key columns}\\
\midrule
\code{investigation} & psi\_id, name, platform, doi, url\\
\code{source\_document} & id, investigation\_id, kind (table / report / image), filename, url, sha256, version\\
\code{experiment} & id, investigation\_id, source\_document\_id, source\_row or page, quote, ten condition columns (each \code{value}, \code{unit}, \code{status}, \code{source}), raw\_row JSONB\\
\code{outcome} & experiment\_id, metric (spread\_rate, burn\_time, extinction, ignited, flame\_length, \dots), value, unit, status, source\\
\code{review} & experiment\_id, reviewer, decision (approved / corrected / rejected), diff JSONB, timestamp\\
\code{model\_version} & name, version, trained\_on (dataset hash), metrics JSONB, domain JSONB, created\_at\\
\code{audit\_log} & every normalization decision and correction --- the auditable ``agentic memory'' from the README\\
\bottomrule
\end{tabularx}

\medskip
Store condition columns \emph{flat} (\code{oxygen\_value}, \code{oxygen\_status}, \code{oxygen\_source})
rather than as one JSON blob. Flat columns make range search and the gap map fast and simple.

% ================= 4 =================
\section{Step-by-step build}

Each step ends with a \textbf{Done when} check. Sizing assumes a team of 2--4 people.

\subsection{Step 1 --- Clean up the repository and set up tooling \hfill\normalfont\small ($\tfrac{1}{2}$ week)}
\begin{enumerate}
  \item Create \code{apps/web}, \code{apps/api}, \code{pipeline}, \code{ml}, \code{schema}; move files with \code{git mv} to keep history.
  \item Web: add TypeScript (\code{allowJs: true}), Vitest, ESLint and Prettier.
  \item API: \code{uv init apps/api}; add \code{fastapi uvicorn sqlalchemy alembic psycopg pydantic pint pandas pandera}.
  \item \code{dvc init}; move \code{psi\_data} and \code{ntrs\_pdfs} into \code{data/raw}, \code{dvc add} it, configure a remote.
  \item Add a \code{LICENSE} (the README notes there is none) and \code{.env.example}.
  \item GitHub Actions: \code{web.yml} (lint, Vitest, build) and \code{api.yml} (ruff, pytest).
\end{enumerate}
\done{\code{npm run dev}, \code{npm test}, \code{uv run pytest} and CI all pass on a clean clone; the existing five science tests stay green.}

\subsection{Step 2 --- Define the schema in code \hfill\normalfont\small ($\tfrac{1}{2}$ week)}
\begin{enumerate}
  \item Write \code{schemas.py} (Pydantic) for \code{Field}, the ten conditions, \code{Outcome} and \code{Experiment}.
  \item Write SQLAlchemy models and the first Alembic migration.
  \item Generate \code{schema/flamme-experiment.schema.json} with \code{Experiment.model\_json\_schema()}.
  \item Generate TypeScript types from the OpenAPI spec with \code{openapi-typescript}.
  \item Port the five curated records in \code{src/data.js} into a seed file. They are the golden test set.
\end{enumerate}
\done{the five seed records round-trip DB $\rightarrow$ API $\rightarrow$ UI unchanged, including \code{037 mm} staying \code{ambiguous} and pressure staying \code{not\_reported}.}

\subsection{Step 3 --- Table ingestion pipeline \hfill\normalfont\small ($1\tfrac{1}{2}$ weeks, deliverable 1)}
PSI experimental tables all have different columns, so write \textbf{one adapter per investigation}, not one generic parser.
\begin{enumerate}
  \item For each of the 22 table files, write \code{pipeline/adapters/<investigation>.py} with a \code{COLUMN\_MAP} (source column $\rightarrow$ condition) and a \code{parse\_row()} that keeps \code{source} verbatim.
  \item \code{normalize.py}: convert units with pint; map vocabularies (``co-current'', ``cocurrent'' $\rightarrow$ \code{concurrent}); set status (\code{\textasciitilde} $\rightarrow$ approximate, words $\rightarrow$ qualitative, unparseable $\rightarrow$ ambiguous, blank $\rightarrow$ not\_reported).
  \item \code{validate.py} (pandera): $0 \le X_{\mathrm{O_2}} \le 100$, $0 \le P \le 200$ kPa, $u \ge 0$, thickness $>0$, and every row has a source file and row number.
  \item \code{load.py}: idempotent upsert keyed on (source document, source row).
  \item Lift the parsing logic already in \code{findings/build\_findings.py} into the adapters.
\end{enumerate}
Priority: the investigations the plan names --- \textbf{Saffire} (PSI-99 etc.), \textbf{BASS / BASS-II} (PSI-25), \textbf{FLEX} (PSI-69) --- then the rest.
\done{\code{make ingest} loads every table investigation; each adapter has a pytest checking row counts and three hand-verified rows; a validation report shows the ``not reported'' rate per field.}

\subsection{Step 4 --- Report extraction with human review \hfill\normalfont\small ($1\tfrac{1}{2}$ weeks, deliverable 1)}
Many conditions (pressure, ignition energy, flow direction) appear only in NTRS reports, drop-tower and parabolic-flight papers, and NASA-STD-6001 test reports.
\begin{enumerate}
  \item \code{extract\_pdf/}: PyMuPDF for text, pdfplumber for tables; chunk by page.
  \item Ask an LLM for \textbf{structured output} matching the Pydantic \code{Experiment} schema, with a \code{page} and verbatim \code{quote} for every non-null field. Reject any field whose quote is not found in the page text (string match in code).
  \item Everything extracted lands as \code{pending}; nothing pending is shown publicly.
  \item Build a \code{/review} screen (reviewer login): PDF page on the left, extracted fields on the right; approve, correct or reject. Corrections go to \code{audit\_log}.
\end{enumerate}
\done{all 36 downloaded NTRS PDFs are processed; at least one reviewer has approved records from each; a test confirms unapproved rows never appear in public API responses.}

\subsection{Step 5 --- Search and compare \hfill\normalfont\small (1 week, deliverable 2)}
\begin{enumerate}
  \item API: \code{GET /experiments?material\_class=PMMA\&o2\_lt=21\&direction=opposed\&flow\_lt=10}; each filter maps to SQL on the flat columns.
  \item Show how unknowns behave: a row with oxygen \code{not\_reported} is excluded from an \Otwo{} filter but counted in a ``12 more records have unknown \Otwo'' link.
  \item Free-text bar: a deterministic parser (regex + vocabulary) turns the plan's example query into a filter object shown back as editable chips. An optional LLM fallback may \emph{propose} filters; the user always sees the chips first.
  \item Keep the \code{science.js} similarity ranking as ``closest experiments'', now server-side over the whole database.
  \item Compare view: up to four experiments as columns, the ten conditions as rows, standard units, unknown cells visibly marked, each cell linked to its source row or page.
  \item Filters live in the URL so searches can be shared.
\end{enumerate}
\done{the plan's example query returns the right rows; a Playwright test covers search $\rightarrow$ compare $\rightarrow$ open source link.}

\subsection{Step 6 --- Gap map \hfill\normalfont\small (1 week, deliverable 3)}
\begin{enumerate}
  \item API \code{GET /gaps?x=oxygen\&y=flow\&facet=gravity} returns counts per bin. Default bins: \Otwo{} 2 percentage points, flow 5 cm/s, pressure 10 kPa.
  \item Heatmap: coloured cells are tested (count shown), outlined empty cells are \textbf{gaps}, hatched cells contain only approximate values.
  \item Clicking a cell opens the search results for that bin. Facets: gravity, flow direction, material class.
  \item Overlay the fitted flammability boundary from Step 7, so users see conditions that are untested \emph{and} near the boundary --- the most useful place for a next experiment.
\end{enumerate}
\done{map counts match a direct SQL query (pytest), and clicking any tested cell returns exactly that many rows.}

\subsection{Step 7 --- Fit models to the data \hfill\normalfont\small (2 weeks, deliverable 4: science)}
Data is small and heterogeneous, so use simple, physics-informed, interpretable models --- one per question, never one model for everything.

\medskip
\begin{tabularx}{\linewidth}{@{}p{2.8cm} p{3.6cm} p{3.4cm} L@{}}
\toprule
\textbf{Model} & \textbf{Inputs} & \textbf{Target} & \textbf{Method}\\
\midrule
Flammability boundary & $p_{\mathrm{O_2}}$, $u$, material class, gravity & ignited and sustained (yes/no) & Logistic regression or Gaussian-process classifier\\
Spread rate & thickness, $p_{\mathrm{O_2}}$, $u$, direction & spread rate (mm/s), only where sources measure the same quantity & Physics-form regression\\
Suppression & agent, concentration, $u$ & extinguished (yes/no), time to extinction & Logistic regression or survival model\\
\bottomrule
\end{tabularx}

\medskip
For spread over a \textbf{thermally thin} sample in opposed flow, use the de~Ris result as the model shape instead of a black box:
\[
  V_f \;=\; \dfrac{\pi}{4}\,\dfrac{\lambda_g\,(T_f - T_v)}{\rho_s\, c_s\, \tau\,(T_v - T_\infty)}
\]
so spread rate should scale as $V_f \propto \dfrac{1}{\tau}$ with sample thickness $\tau$. Fit the coefficient (and any extra oxygen and flow dependence) from data, and report how well the observations follow this shape.

\medskip\noindent\textbf{Rules for every fit}
\begin{enumerate}
  \item \textbf{Grouped cross-validation by investigation} (\code{GroupKFold}) so one experiment never sits in both train and test.
  \item Start with a baseline (mean or single variable); keep a model only if it beats it.
  \item Report calibrated uncertainty (prediction intervals or calibrated probabilities).
  \item Record the \textbf{domain} (ranges or convex hull of training inputs). Outside it the UI says ``outside tested conditions'' instead of predicting.
  \item \code{ml/export.py} writes \code{models/flammability@1.0.json} (coefficients, domain, metrics, dataset hash) and registers it in \code{model\_version}.
\end{enumerate}
\done{each model card shows cross-validated metrics against the baseline, the domain and the dataset version; a pytest confirms export $\rightarrow$ reload gives identical predictions.}

\subsection{Step 8 --- Ground the 3D lab in data \hfill\normalfont\small ($1\tfrac{1}{2}$ weeks, deliverable 4: visual)}
\begin{enumerate}
  \item Split \code{derive(p)} in \code{labModel.js} into two layers:
  \begin{itemize}
    \item \code{predict(p)}: physical quantities from the exported model JSON (ignition probability, spread rate with interval, extinguished or not);
    \item \code{toVisual(prediction)}: flame size, soot, shell shape and colour, keeping the existing Three.js look.
  \end{itemize}
  \item Load models with \code{fetch('/models/flammability@1.0.json')}; logistic and linear models run in a few lines of JavaScript in the browser.
  \item Out of domain: keep the illustrative curve and show the banner \emph{``Illustrative --- no NASA data for these conditions''}.
  \item Add a \textbf{Closest real experiments} panel: query the similarity endpoint with the current slider values and show the top three with real outcome and source link.
  \item Label every on-screen number \textbf{Observed}, \textbf{Predicted (model@version)} or \textbf{Unknown}.
\end{enumerate}
\done{for each seed experiment, setting the sliders to its conditions shows it as the closest experiment, and the predicted outcome direction agrees with the observed one --- or the mismatch is shown, not hidden.}

\subsection{Step 9 --- Reporting template \hfill\normalfont\small ($\tfrac{1}{2}$ week, deliverable 5)}
\begin{enumerate}
  \item Publish \code{schema/flamme-experiment.schema.json} (generated in Step 2) at \code{/schema}.
  \item Generate a fillable CSV template and a web form from the same schema (\code{react-jsonschema-form}) that flags missing conditions.
  \item A ``Validate my file'' page: upload CSV or JSON, receive a completeness report for the ten conditions.
\end{enumerate}
\done{a file exported from the database validates against the schema; a file missing pressure gets a clear warning.}

\subsection{Step 10 --- Outreach mode \hfill\normalfont\small ($\tfrac{1}{2}$ week, deliverable 6)}
\begin{enumerate}
  \item Keep the 92-second \code{/\#studio-play} sweep; captions now cite real experiments (``Saffire II sample 2-5 at about 22\% \Otwo{} \dots'').
  \item A classroom route with large text and no admin UI, available offline as a PWA (\code{vite-plugin-pwa}).
\end{enumerate}
\done{the sweep runs end to end on a school laptop in Chrome or Firefox at 30 fps or more.}

\subsection{Step 11 --- Quality, accessibility and deployment \hfill\normalfont\small (1 week, overlapping)}
\begin{enumerate}
  \item Tests: Vitest for \code{labModel} and \code{science}; Playwright for search $\rightarrow$ compare $\rightarrow$ lab $\rightarrow$ source; pytest for adapters, API and model export; a visual snapshot of the lab.
  \item Accessibility: keep the keyboard shortcuts, add \code{aria-live} for lab status, audit with axe.
  \item Performance: lazy-load the Three.js lab chunk; cap device pixel ratio at 2.
  \item Deploy: web to Vercel (\code{VITE\_API\_URL}); API to Render or Fly (\code{DATABASE\_URL}); database on Neon; a release job runs \code{alembic upgrade head} and \code{make ingest}.
  \item Footer on every page: ``Not a NASA product. Not for certification or operational safety decisions.''
\end{enumerate}

% ================= 5 =================
\section{Timeline}

\begin{tabularx}{\linewidth}{@{}p{2cm} p{3.2cm} L@{}}
\toprule
\textbf{Week} & \textbf{Steps} & \textbf{Milestone}\\
\midrule
1 & 1, 2 & Tooling, schema, five seed records through the full stack\\
2--3 & 3 & All PSI tables in the database\\
3--4 & 4, 5 & Reports extracted and reviewed; search and compare live\\
5 & 6, 9 & Gap map; reporting template\\
6--7 & 7 & Fitted models with model cards\\
7--8 & 8, 10, 11 & Data-grounded lab, outreach mode, deployment\\
\bottomrule
\end{tabularx}

\begin{keybox}
\textbf{48-hour hackathon path.} Steps 1--3 for Saffire, BASS and FLEX only; Steps 5 and 6 with DuckDB-WASM over
exported JSON; one flammability-boundary model for Steps 7--8; and the closest-experiments panel. That alone
demonstrates the full idea end to end with real data.
\end{keybox}

% ================= 6 =================
\section{Suggested team split}

\begin{tabularx}{\linewidth}{@{}p{3.2cm} L@{}}
\toprule
\textbf{Role} & \textbf{Owns}\\
\midrule
Data & Steps 3 and 4: adapters, normalization, PDF extraction, review\\
Backend & Step 2, the APIs for Steps 5 and 6, deployment\\
Frontend / 3D & The interfaces for Steps 5 and 6, plus Steps 8 and 10\\
ML / science & Step 7, model cards, checking physics and limitations\\
\bottomrule
\end{tabularx}

% ================= 7 =================
\section{Risks and mitigations}

\begin{tabularx}{\linewidth}{@{}p{5cm} L@{}}
\toprule
\textbf{Risk} & \textbf{Mitigation}\\
\midrule
Too few comparable outcome measurements for a model & Fit per question and per material class; if data is too thin, show only nearest experiments with ``insufficient data''\\
LLM extraction invents values & Page and quote required, quote string-matched in code, human approval before publishing\\
Unit or definition mismatch (e.g.\ Saffire ``Spread Length'' column holding mm/s) & Keep \code{source} verbatim, mark \code{ambiguous}, exclude from fits until reviewed\\
Leakage inflating model scores & \code{GroupKFold} by investigation\\
Users treat the lab as a real fire prediction & Observed / Predicted / Unknown labels, out-of-domain banner, disclaimer footer\\
Raw data size (about 14 TB on PSI) & Tables and reports only; raw downloads stay opt-in and size-capped\\
\bottomrule
\end{tabularx}

\end{document}
